\chapter{度量空间}
\section{基础空间}
\begin{definition}[距离空间、度量Metric Space]
所谓度量空间，就是指对偶 $(X,d)$，其中 $X$ 是一个集合，
$d$ 是 $X$ 上的一个度量（或 $X$ 上的距离函数），即 $d$ 是定义在
$X\times X$ 上的函数，并且对所有 $x,y,z\in X$ 满足以下四条公理：
\begin{align*}
(\mathrm{M}_1) \quad & d(x,y) \text{ 是实值、有限且非负的}; \\
(\mathrm{M}_2) \quad & d(x,y)=0 \iff x=y; \\
(\mathrm{M}_3) \quad & d(x,y)=d(y,x) \quad \text{(对称性)}; \\
(\mathrm{M}_4) \quad & d(x,y)\le d(x,z)+d(z,y) \quad \text{(三角不等式)}.
\end{align*}
\end{definition}
为叙述方便，我们给出下面几个有关的术语。$X$ 叫作 $(X,d)$ 的基集，$X$ 的元素叫作空间 $(X,d)$ 的点。给定 $x,y \in X$，我们把非负实数 $d(x,y)$ 叫作 $x$ 和 $y$ 之间的距离。M$_1$ 到 M$_4$ 叫作度量公理。“三角不等式”的名字是受初等几何的启发而得到的，如图~1-2 所示。

\begin{center}
\begin{tikzpicture}[scale=2]
  % 点
  \coordinate (x) at (0,0);
  \coordinate (z) at (2,0);
  \coordinate (y) at (2.5,1.2);

  % 三角形边
  \draw (x) -- (y) -- (z) -- cycle;

  % 标注
  \node[below left] at (x) {$x$};
  \node[below right] at (z) {$z$};
  \node[above] at (y) {$y$};

  \node[above left] at ($(x)!0.5!(y)$) {$d(x,y)$};
  \node[below] at ($(x)!0.5!(z)$) {$d(x,z)$};
  \node[right] at ($(z)!0.5!(y)$) {$d(z,y)$};
\end{tikzpicture}

图 1-2 \quad 平面上的三角不等式
\end{center}
用归纳法可以从 $M_4$ 推得广义三角不等式
\[
d(x_1,x_n) \le d(x_1,x_2) + d(x_2,x_3) + \cdots + d(x_{n-1},x_n), \tag{1.1.1}
\]

在不引起混淆的情况下，我们常将度量空间 $(X,d)$ 简写成 $X$。  
如果取子集 $Y \subseteq X$ 且把 $d$ 限制在 $Y \times Y$ 上，则可得 $(Y,\tilde d)$，因此 $Y$ 上的度量就是限制：
\[
\tilde d = d|_{Y \times Y}.
\]

$\tilde d$ 叫作 $d$ 在 $Y$ 上导出的度量。
\begin{definition}[欧几里得空间 $\mathbb{R}^n$、酉空间 $\mathbb{C}^n$、复平面 $\mathbb{C}$的距离]
前面几个例子是 $n$ 维欧几里得空间 $\mathbb{R}^n$ 的特殊情况。
如果取所有形如 $x=(\xi_1,\dots,\xi_n),\, y=(\eta_1,\dots,\eta_n)$ 的
$n$ 元实序列集合作为基集，用
\[
d(x,y) = \sqrt{(\xi_1-\eta_1)^2 + \cdots + (\xi_n-\eta_n)^2}\ (\ge 0) \tag{1.1.6}
\]
在其上定义欧几里得度量，便得到这个空间。

$n$ 维酉空间 $\mathbb{C}^n$ 的基集是所有 $n$ 元复序列的集合，其度量定义为
\[
d(x,y) = \sqrt{|\xi_1-\eta_1|^2 + \cdots + |\xi_n-\eta_n|^2}\ (\ge 0). \tag{1.1.7}
\]

当 $n=1$ 时，便得到复平面 $\mathbb{C}$，其普通度量为
\[
d(x,y) = |x-y|. \tag{1.1.8}
\]

（$\mathbb{C}^n$ 有时又叫作 $n$ 维复欧几里得空间。）
\end{definition}
\begin{definition}[一致有界序列空间$l^\infty$]
这个例子和下一个例子给人的第一个印象是，度量空间的概念具有如此惊人的一般性。
我们取所有有界的复数序列的集合作为基集 $X$；也就是说，$X$ 中的每个元素都是形如
\[
x = (\xi_1,\xi_2,\dots) \quad \text{简记为 } x=(\xi_j)
\]
的复数序列，且对于 $j=1,2,\dots$ 有
\[
|\xi_j|\le c_x,
\]
其中 $c_x$ 是依赖于 $x$ 的实数，但与 $j$ 无关。

我们在 $X$ 上用
\[
d(x,y) = \sup_{j\in \mathbb{N}} |\xi_j - \eta_j| \tag{1.1.9}
\]
来定义度量，其中 $y=(\eta_j)\in X$ 且 $\mathbb{N}=\{1,2,\dots\}$，
而 $\sup$ 表示上确界（最小上界）。

这样得到的度量空间通常记为 $l^\infty$。
（这个有点奇怪的记法由 1.2-3 导出。因为 $X$ 中的每个元素（或 $X$ 的每个点）都是序列，所以 $l^\infty$ 是一个序列空间。）
\end{definition}
每一个元素的模长都是有上界的。$d(x,y)$ 本质上就是 两个序列在所有位置上的最大差距。
\begin{definition}[1.2-1 序列空间 $s$]
这个空间的基集 $X$ 是所有（有界或无界的）复数序列的集合，其度量 $d$ 定义为
\[
d(x,y) = \sum_{j=1}^{\infty} \frac{1}{2^j} \cdot \frac{|\xi_j - \eta_j|}{1 + |\xi_j - \eta_j|},
\]
其中 $x = (\xi_j), \; y = (\eta_j)$。
注意，1.1-6 中的度量在这里是不合适的。（为什么？）
\end{definition}
\begin{proof}
容易看出，公理 $M_1 \sim M_3$ 是满足的，让我们来验证 $M_4$。为此，我们利用定义在 $\mathbb{R}$ 上的辅助函数
\[
f(t) = \frac{t}{1+t}.
\]

对它微分可得
\[
f'(t) = \frac{1}{(1+t)^2},
\]
显然对任意 $t \in \mathbb{R}$ 有 $f'(t) > 0$。因此 $f$ 是单调递增的。从而由
\[
|a+b| \leq |a| + |b|
\]
可推出
\[
f(|a+b|) \leq f(|a| + |b|).
\]

将上述函数写出并应用关于数的三角不等式，便得到
\[
\frac{|a+b|}{1+|a+b|}
\leq \frac{|a|+|b|}{1+|a|+|b|}
= \frac{|a|}{1+|a|+|b|} + \frac{|b|}{1+|a|+|b|}
\leq \frac{|a|}{1+|a|} + \frac{|b|}{1+|b|}.
\]
\end{proof}
\begin{definition}[S空间]

空间 $S$。

与例 1.1.3 类似，设 $E\subset\mathbb{R}$ 是一个 Lebesgue 可测集，且 $0<m(E)<\infty$。
考虑 $E$ 上几乎处处有穷的可测函数全体（其中几乎处处相等的函数看成是同一元）。
定义
\[
d(x,y) \;=\; \int_{E} \frac{|x(t)-y(t)|}{1+|x(t)-y(t)|}\,dt .
\]
与例 1.1.3 的证明类似，这是一个距离空间，把这个空间记为 $S$。
   
\end{definition}
\begin{definition}[函数空间 $C[a,b]$]
设 $J=[a,b]$ 是一个闭区间。我们考虑所有定义在 $J$ 上的连续实值函数
\[
x(t), \; y(t), \;\dots
\]
的集合作为基集 $X$。

在 $X$ 上定义度量
\[
d(x,y) = \max_{t \in J} |x(t) - y(t)|.
\tag{1.1.10}
\]

这里 $\max$ 表示最大值。由此得到的度量空间记作 $C[a,b]$，其中字母 $C$ 来自英文单词 \emph{Continuous} 的第一个字母。因为 $C[a,b]$ 的每个元素都是一个函数，所以 $C[a,b]$ 是一个函数空间。
\end{definition}基集：就是“讨论的内容”（对象集合）。度量讨论的标准。要求定义在J上面的连续实值函数。度量为两个函数距离最远处。
\begin{definition}[离散度量空间]
设 $X$ 是任意一个集合，定义映射 $d:X\times X \to \mathbb{R}$ 为
\[
d(x,y) =
\begin{cases}
0, & x = y, \\
1, & x \neq y.
\end{cases}
\]
则 $(X,d)$ 称为离散度量空间。
\end{definition}
\begin{definition}[$l^p$ 空间]
设 $p \geq 1$ 为一个固定的实数。  
称所有实数（或复数）序列 $x=(\xi_1,\xi_2,\dots)$，若满足
\begin{equation}
\sum_{j=1}^{\infty} |\xi_j|^p < \infty, 
\tag{1.2.1}
\end{equation}
则称 $x$ 属于 $l^p$ 空间。

在 $l^p$ 空间上，定义度量
\begin{equation}
d(x,y) = \left( \sum_{j=1}^{\infty} |\xi_j - \eta_j|^p \right)^{1/p},
\quad x=(\xi_j),~ y=(\eta_j).
\tag{1.2.2}
\end{equation}

若取所有满足 {1.2.1} 的实数序列，得到实 $l^p$ 空间，记作 $l^p_{\mathbb{R}}$；  
若取所有满足 {1.2.1} 的复数序列，得到复 $l^p$ 空间，记作 $l^p_{\mathbb{C}}$。

特别地，当 $p=2$ 时，$l^2$ 空间是一个希尔伯特空间。
\end{definition}
\begin{definition}[共轭指数]
令 $p>1$，并定义 $q$ 满足
\begin{equation}
\frac{1}{p} + \frac{1}{q} = 1. \tag{1.2.4}
\end{equation}
则称 $p$ 和 $q$ 为一对 \textbf{共轭指数}。

由 {1.2.4} 可得
\begin{equation}
1 = \frac{p+q}{pq}, \quad pq = p+q, \quad (p-1)(q-1)=1. \tag{1.2.5}
\end{equation}
\end{definition}

因此 
\[
\frac{1}{p-1} = q-1.
\]
若令 $u = t^{p-1}$，则 $t = u^{\,\frac{1}{p-1}} = u^{\,q-1}$。

设 $\alpha,\beta$ 为任意两个正数，由于 $\alpha\beta$ 可由图形面积表示，通过积分可得不等式
\begin{equation}
\alpha \beta \leq \int_{0}^{\alpha} t^{p-1} \, dt + \int_{0}^{\beta} u^{q-1} \, du 
= \frac{\alpha^p}{p} + \frac{\beta^q}{q}. \tag{1.2.6}
\end{equation}

注意，当 $\alpha=0$ 或 $\beta=0$ 时，上述不等式依然成立。
% 归一化条件
\begin{equation}
\sum_{j=1}^{\infty} |\tilde\xi_j|^{\,p} = 1,
\qquad
\sum_{j=1}^{\infty} |\tilde\eta_j|^{\,q} = 1.
\tag{1.2.7}
\end{equation}

% 由 Young 不等式得到逐项不等式
令 $\alpha = |\tilde\xi_j|$、$\beta = |\tilde\eta_j|$，将 \,(1.2.6)\, 代入得
\[
|\tilde\xi_j \tilde\eta_j|
\le \frac{1}{p}\,|\tilde\xi_j|^{\,p} + \frac{1}{q}\,|\tilde\eta_j|^{\,q}.
\]

两端对 $j$ 求和，并利用 \,(1.2.7)\, 及 \,(1.2.4)\,（即 $1/p + 1/q = 1$），便得到
\begin{equation}
\sum_{j=1}^{\infty} \bigl|\tilde\xi_j \tilde\eta_j\bigr|
\le \frac{1}{p}\sum_{j=1}^{\infty} |\tilde\xi_j|^{\,p}
   + \frac{1}{q}\sum_{j=1}^{\infty} |\tilde\eta_j|^{\,q}
= \frac{1}{p} + \frac{1}{q} = 1.
\tag{1.2.8}
\end{equation}

% 任意非零序列的归一化
现在取任意非零序列 $x=(\xi_j)\in \ell^{p}$ 与 $y=(\eta_j)\in \ell^{q}$，并置
\begin{equation}
\tilde\xi_j = \frac{\xi_j}{\left(\displaystyle\sum_{k=1}^{\infty}|\xi_k|^{\,p}\right)^{1/p}},
\qquad
\tilde\eta_j = \frac{\eta_j}{\left(\displaystyle\sum_{m=1}^{\infty}|\eta_m|^{\,q}\right)^{1/q}},
\tag{1.2.9}
\end{equation}
则由定义可见它们满足 \,(1.2.7)\,。
由它们满足 (1.2.7)，所以能够应用不等式 (1.2.8)。
把 (1.2.9) 代入 (1.2.8)，再用 (1.2.9) 的分母的乘积去乘所得到不等式的两端，
便得到关于和式的 
\begin{theorem}{Hölder赫尔德 不等式}

\begin{equation}
\sum_{j=1}^{\infty} |\xi_j \eta_j|
\le
\left( \sum_{k=1}^{\infty} |\xi_k|^p \right)^{1/p}
\left( \sum_{m=1}^{\infty} |\eta_m|^q \right)^{1/q},
\qquad (p>1,\; \tfrac{1}{p}+\tfrac{1}{q}=1).
\tag{1.2.10}
\end{equation}
\end{theorem}
这个不等式来源于 Hölder (O.~Hölder, 1889)。

特别地，当 $p=2,\; q=2$ 时，(1.2.10) 就给出关于和式的
\begin{theorem}{柯西施瓦茨Cauchy–Schwarz 不等式}：

\begin{equation}
\sum_{j=1}^{\infty} |\xi_j \eta_j|
\le
\sqrt{ \sum_{k=1}^{\infty} |\xi_k|^2 } \;
\sqrt{ \sum_{m=1}^{\infty} |\eta_m|^2 }.
\tag{1.2.11}
\end{equation}
\end{theorem}
\begin{theorem}[闵可夫斯基不等式]
    \begin{equation}
\Biggl(\sum_{j=1}^{\infty} \bigl|\xi_j + \eta_j \bigr|^p \Biggr)^{1/p}
\le
\Biggl(\sum_{k=1}^{\infty} |\xi_k|^p \Biggr)^{1/p}
+
\Biggl(\sum_{m=1}^{\infty} |\eta_m|^p \Biggr)^{1/p}.
\tag{1.2.12}
\end{equation}
\end{theorem}


\section{课本1.1}
\begin{definition}[收敛与极限]
设 $\{x_n\}$ 是距离空间 $(X,d)$ 中的一个点列，$x_0$ 是 $X$ 中一点，  
如果当 $n \to \infty$ 时，$d(x_n, x_0) \to 0$，则称当 $n \to \infty$ 时，$\{x_n\}$ 以 $x_0$ 为极限，  
或当 $n \to \infty$ 时，$\{x_n\}$ 收敛于 $x_0$。记为
\[
x_n \to x_0 \quad (n \to \infty),
\]
或
\[
\lim_{n \to \infty} x_n = x_0.
\]
\end{definition}
\begin{theorem}[极限唯一，子列收敛到同一个点]
设 $\{x_n\}$ 是距离空间 $X$ 中的收敛点列，则：
\begin{enumerate}
  \item $\{x_n\}$ 的极限是唯一的；
  \item 如果 $x_0$ 是 $\{x_n\}$ 的极限，那么 $\{x_n\}$ 的任一子列 $\{x_{n_k}\}$ 必收敛且以 $x_0$ 为极限。
\end{enumerate}
\end{theorem}
\begin{theorem}[度量连续性metric continuity]
设 $(X,d)$ 是距离空间，则
\[
\big| d(x,y) - d(x_1,y_1) \big|
\leq d(x,x_1) + d(y,y_1), 
\qquad (x,y,x_1,y_1 \in X).
\]
\end{theorem}
三角不等式的推广。\\
\noindent\textbf{三、距离空间的连续映射、等距}

设 $(X,d), (X_1,d_1)$ 是距离空间，$f: X \to X_1$ 是一个映射，$x_0 \in X$，如果  
$\forall \varepsilon > 0$，存在 $\delta > 0$，使得 $d(x,x_0) \leq \delta$ 的一切 $x \in X$，  
\[
d_1\big(f(x), f(x_0)\big) < \varepsilon,
\]
则称映射 $f$ 在 $x_0$ 点连续；如果 $f$ 在 $X$ 上的每一点连续，则称 $f$ 在 $X$ 上连续。  

从以上定义可以看出，距离空间间的连续映射是分析中大家熟知的连续函数的推广。  
\begin{definition}[等距映射]
    

如果对任意 $x,y \in X$，  
\[
d_1(f(x), f(y)) = d(x,y),
\]
则称 $f$ 是一个等距映射。  
\end{definition}
一个等距映射一定是一个连续映射，并且是一个一一映射。  

对于两个距离空间 $(X,d), (X_1,d_1)$，如果存在一个映上的等距映射 $f: X \to X_1$，则称 $(X,d)$ 与 $(X_1,d_1)$ 等距。  

从距离空间的观点来看，两个等距的距离空间没有本质的差别，因此，今后我们把两个等距的距离空间看成是同一空间。
\begin{theorem}[距离空间性质]
设 $X$ 是距离空间，则 $X$ 中的开集具有以下基本性质：
\begin{enumerate}
  \item 全空间 $X$ 与空集 $\varnothing$ 是开集；
  \item 任意个开集的并集是开集；
  \item 任意有限多个开集的交集是开集。
\end{enumerate}
\end{theorem}
\section{距离空间的点集}
\begin{proposition}[内部，接触点，闭包，极限点]
    

设 $A$ 是距离空间中的任一集，用 $A^{\circ}$ 表示所有 $A$ 的开子集的并集，
称 $A^{\circ}$ 为集 $A$ 的\textbf{内部}。由定理 1.3，$A^{\circ}$ 是开集，
并且是包含在 $A$ 中的最大的开集。

设 $A$ 是距离空间 $X$ 中的点集，$x_0 \in X$，如果 $\forall \varepsilon >0$，
球 $S(x_0,\varepsilon)$ 中都包含 $A$ 的点，即
\[
S(x_0,\varepsilon) \cap A \neq \varnothing,
\]
称 $x_0$ 为集 $A$ 的\textbf{接触点}(聚点)，集 $A$ 的接触点的全体称为 $A$ 的\textbf{闭包}，
$A$ 的闭包记为 $\overline{A}$，显然 $A \subset \overline{A}$。
但是，$A$ 的接触点不必属于 $A$。

如果 $\forall \varepsilon > 0$，球 $S(x_0,\varepsilon)$ 中总包含 $A$ 中
不同于 $x_0$ 的点，称 $x_0$ 是 $A$ 的\textbf{极限点}。
显然，$A$ 的极限点必是 $A$ 的接触点，反之则不然。
\end{proposition}
\begin{proposition}[闭集的刻画]
设 $E$ 是一个集合，则以下条件等价：
\begin{enumerate}
  \item 若任意 $\{x_n\} \subset E$，当 $\lim\limits_{n \to \infty} x_n = x$ 时，有 $x \in E$，则称 $E$ 是闭集；
  \item 若 $E = \overline{E}$，则 $E$ 是闭集；
  \item $\overline{E} = E \cup \partial E = \mathring{E} \cup \partial E = \{E \text{ 的全部聚点}\} \cup \{E \text{ 的全部孤立点}\}$；
  \item 任何有限点集均是闭集。
\end{enumerate}
\end{proposition}
\begin{theorem}[集A和闭包的关系]
设 $A,B$ 是距离空间 $X$ 的子集，则：
\begin{enumerate}
  \item $A \subset \overline{A}$；
  \item $\overline{\overline{A}} = \overline{A}$；
  \item $\overline{A \cup B} = \overline{A} \cup \overline{B}$；
  \item $\overline{\varnothing} = \varnothing$。
\end{enumerate}
\end{theorem}
设 $A$ 是距离空间 $X$ 中的集，如果 $A = \overline{A}$，则称 $A$ 为\textbf{闭集}。

由以上定理，任一集的闭包 $\overline{A}$ 是闭集，它是包含集 $A$ 的最小闭集。

不难证明，在一个距离空间中，任一开集的余集是闭集，任一闭集的余集是开集。
因此，由 de Morgan 公式立刻可得闭集的基本性质。
\begin{theorem}[闭集性质]
设 $X$ 是距离空间，则：
\begin{enumerate}
  \item 全空间 $X$ 及空集 $\varnothing$ 是闭集；
  \item 任意个闭集的交集是闭集；
  \item 任意有限个闭集的并集是闭集。
\end{enumerate}
\end{theorem}
\begin{definition}[稠密子集]
设 $X$ 是一个度量空间（或拓扑空间），$A \subseteq X$。  
如果对任意 $x \in X$ 和任意 $\varepsilon > 0$，存在 $a \in A$ 使得
\[
d(x,a) < \varepsilon,
\]
那么称 $A$ 在 $X$ 中是 \textbf{稠密的 (dense in $X$)}。

等价地说：
\[
\overline{A} = X,
\]
即 $A$ 的闭包等于整个空间 $X$。
\end{definition}

\begin{remark}[直观理解]
\begin{enumerate}
  \item 任意集合均有稠密子集。
  \item 若 $A \subset \overline{B}$，则有
  \[
    \overline{B} \subset \bigcup_{x \in B} S(x,\varepsilon), \quad \forall \varepsilon > 0.
  \]
  \item $\forall a \in A, \; \forall \varepsilon > 0, \; \exists b \in B, \; s.t.\; a \in S(x,\varepsilon)$。
  \item 稠密具有传递性。若 $H$ 在 $A$ 中稠密，$A$ 在 $D$ 中稠密，则 $H$ 在 $D$ 中稠密。
\end{enumerate}
\begin{itemize}
  \item 稠密 = “没有遗漏”：$A$ 里的点可以无限逼近 $X$ 的每一个点。
  \item 稠密子集可能“稀疏”，但它覆盖了空间的“整体形状”。
\end{itemize}
\end{remark}
\begin{definition}[可分性 separable。]
定义1.2.1:设 $X$ 是距离空间，如果 $X$ 中存在一个稠密可数子集，则称 $X$ 是可分的。  

对于子集 $A \subset X$，如果 $X$ 中存在可数子集 $B$，使得 $B$ 在 $A$ 中稠密，则称集 $A$ 是可分的。
\end{definition}

\begin{remark}[词源解释]
“Separable” 来自 \textit{separate}（分开、区分）。  
含义是：虽然空间可能非常“大”，但它能被一个可数集“分辨/刻画”出来。  
换句话说：通过那个稠密的可数子集，可以“逼近”或“区分”空间里的任意点。
\end{remark}

\begin{example}
\begin{itemize}
  \item 实数集 $\mathbb{R}$ 是可分的，因为有理数集 $\mathbb{Q}$ 是稠密且可数的。
  \item 任意实数点都可以被有理数“无限逼近”，所以 $\mathbb{Q}$ 足够“分辨”整个 $\mathbb{R}$。
\end{itemize}
\end{example}

\section{完备距离空间}
我们希望描述一种情况：当序列“走得很远”以后，它的所有项彼此之间越来越接近。
\begin{definition}[Cauchy 列与完备性]
设 $\{x_n\}$ 是距离空间 $X$ 中的点列，如果对任意的 $\varepsilon > 0$，存在自然数 $N$，当 $m,n > N$ 时，
\[
d(x_m, x_n) < \varepsilon,
\]
称 $\{x_n\}$ 是一个 \textbf{Cauchy 列}（基本列）。如果 $X$ 中任意 Cauchy 列都收敛，称距离空间 $X$ 是 \textbf{完备的}。
\end{definition}

由上定义可得以下结论：
\begin{enumerate}
  \item 距离空间中任一收敛点列是 Cauchy 列；
  \item 完备距离空间的任一闭子空间也是完备的。

\begin{example}
在不完备空间中，柯西列可能不收敛。

例如，在有理数集 $\mathbb{Q}$ 中，考虑 $\sqrt{2}$ 的有理数近似列：
\[
1,\, 1.4,\, 1.41,\, 1.414,\, \dots
\]
在有理数的度量下，这个数列是一个柯西列，因为随着 $n$ 的增加，数列的后项之间距离越来越小。

但是，这个数列在 $\mathbb{Q}$ 中没有收敛极限，因为 $\sqrt{2}\notin\mathbb{Q}$。
换句话说，该数列在 $\mathbb{R}$ 中才收敛，其极限为 $\sqrt{2}$。
\end{example}
\end{enumerate}
\begin{dxtips}
    1.完备需要整两个点一个是存在
    如果对任意的 $\varepsilon > 0$，存在自然数 $N$，当 $m,n > N$ 时，
$
d(x_m, x_n) < \varepsilon,
$   \\2.还有他得收敛到具体的一个点
\end{dxtips}
\begin{theorem}[闭球套定理]
设 $X$ 是完备的距离空间，$\overline{S}_n = \overline{S}(x_n, r_n)$ ($n=1,2,\dots$) 是 $X$ 中一列闭球套：
\[
\overline{S}_1 \supset \overline{S}_2 \supset \cdots \supset \overline{S}_n \supset \cdots
\]
且半径 $r_n \to 0 \; (n \to \infty)$，则存在唯一的点 $x \in \bigcap_{n=1}^{\infty} \overline{S}_n$。
\end{theorem}
闭区间套的推广，把闭区间弄成了闭球。
\begin{definition}[疏集与第一/第二纲集]
设 $E$ 是度量空间 $X$ 中的点集。  
如果 $E$ 不在 $X$ 的任何开集中稠密，则称 $E$ 是\textbf{疏集}。  

如果一个集合可以表示为可数个疏集的并集，称为\textbf{第一纲 (meager set) 集}；  
否则称为\textbf{第二纲集}。
\end{definition}
集合 $E$ 为疏集，当且仅当 $\overline{E}$ 不能充满 $X$ 中的任何一个开集，
即
$
\operatorname{int}(\overline{E}) = \varnothing.
$int是 interior（内点集/内部） 的简写。

意思是：集合 A 的内部，也就是包含在 A 中的最大开集。
\begin{example}
设 $X = \mathbb{R}$，考虑整数集 $E = \mathbb{Z}$。

首先，闭包：
\[
\overline{\mathbb{Z}} = \mathbb{Z},
\]
因为整数集中的每个点都是孤立点，没有新的极限点。

其次，内部：
对于任意 $n \in \mathbb{Z}$，若取开区间 $(n-\epsilon, n+\epsilon)$，其中必然含有非整数点
（例如 $n+0.1 \notin \mathbb{Z}$）。
因此不论 $\epsilon > 0$ 多么小，该区间都不可能完全包含在 $\mathbb{Z}$ 中。
所以
\[
\operatorname{int}(\mathbb{Z}) = \varnothing.
\]

由此可见，整数集 $\mathbb{Z}$ 在 $\mathbb{R}$ 中是一个疏集。
直观理解是：整数在实数直线上过于“稀疏”，不能填满任何一段开区间。
\end{example}
\begin{theorem}[Baire 定理]
完备距离空间是第二纲集。
\end{theorem}
\textcolor{blue}{Baire 定理} 告诉我们：如果 $X$ 是一个完备度量空间（比如 $\mathbb{R}, \mathbb{R}^n, \ell^2, C[0,1]$ 等），那么 $X$ 不可能被若干个 \textcolor{blue}{稀疏集合} 盖住。换句话说，完备度量空间是“庞大”的，它的结构保证了空间里一定存在“厚实”的部分。

证明思路是这样的：设想你要用一系列疏集来覆盖整个 $X$，即
\[
X = \bigcup_{n=1}^\infty E_n, \quad \text{其中每个 $E_n$ 是疏集}.
\]
那么，每个 $E_n$ 的闭包 $\overline{E_n}$ 都没有内部点。Baire 定理通过构造递减的闭球序列，保证在完备度量空间里，最终会落到一个非空开集，而这个开集应该落在某个 $\overline{E_n}$ 里，这就和“它没有内部”矛盾。

于是推出：完备度量空间不是第一纲集，也就是说它必然是 \textcolor{blue}{第二纲集}。
\begin{definition}[等距同构]
设 $(X,d)$ 与 $(\widetilde{X},\widetilde{d})$ 是两个度量空间。称 $(X,d)$ 与 $(\widetilde{X},\widetilde{d})$ \textbf{等距同构}，若存在映射 
\[
T: X \to \widetilde{X}
\]
满足：
\begin{enumerate}
    \item $T$ 是满射，即对任意 $y \in \widetilde{X}$，存在 $x \in X$，使得 $Tx = y$；
    \item 对任意 $x,y \in X$，有
    \[
    d(x,y) = \widetilde{d}(Tx,Ty)。
    \]
\end{enumerate}
\end{definition}
第二个条件含了单射因为x≠y，Tx≠Ty
\begin{theorem}[距离空间的完备化定理]
对于任意距离空间 $(X,d)$，必存在一个完备距离空间 $(\tilde{X}, \tilde{d})$，使 $X$ 与 $\tilde{X}$ 的某个稠密子空间 $\tilde{X}_0$ 等距同构，并且 $(\tilde{X}, \tilde{d})$ 在等距同构意义下是唯一的。即若 $(\overline{X}, \overline{d})$ 也是一个完备距离空间，且 $X$ 与 $\overline{X}$ 的某个稠密子空间等距同构，则 $(\tilde{X}, \tilde{d})$ 与 $(\overline{X}, \overline{d})$ 等距同构。
\end{theorem}
多项式空间就不完备他的极限可以到sinx但是里面没有sinx。
{完备化与等价类的构造}  
设 $(X,d)$ 是一个不完备的度量空间。它的 \emph{完备化} $\tilde{X}$ 可以通过柯西列来构造：

\begin{enumerate}
  \item 在 $X$ 中取所有柯西列。  
  \item 定义两个柯西列 $(x_n), (y_n)$ 等价，若
  \[
  (x_n) \sim (y_n) \;\;\Longleftrightarrow\;\; \lim_{n \to \infty} d(x_n,y_n) = 0,
  \]
  即它们“趋向同一个极限点”。  
  \item 每个等价类 $[(x_n)]$ 代表 $\tilde{X}$ 中的一个点：
  \begin{itemize}
    \item 若极限点原本就在 $X$（如常值列 $x_n \equiv a$），则 $[(x_n)]$ 对应于旧点 $a$；
    \item 若极限点不在 $X$（如 $\mathbb{Q}$ 中逼近 $\sqrt{2}$ 的柯西列），则 $[(x_n)]$ 对应于一个“新点”。  
  \end{itemize}
\end{enumerate}

因此，完备化空间 $\tilde{X}$ 就是所有柯西列等价类的集合，并配上自然的度量
\[
\tilde{d}\big( [(x_n)], [(y_n)] \big) := \lim_{n \to \infty} d(x_n,y_n).
\]

\textcolor{blue}{直观理解：等价类 $\approx$ 极限点的替身。完备化就是通过等价类补上原空间缺失的点，从而保证所有柯西列都收敛。}
\begin{example}
    设 $(X,d)$ 是一个距离空间。它的完备化空间记作 $(\tilde{X}, \tilde{d})$，  
其中 $\tilde{X}$ 由 $X$ 中所有柯西列的等价类构成。  

\[
X \;\;\subset\;\; \tilde{X}, \quad 
x \in X \;\longmapsto\; (x_n) \equiv x, \quad x_n = x.
\]

因此，$X$ 在 $\tilde{X}$ 中是稠密子集。
(x)就是全是x的子列，作为代表元。
\end{example}
\begin{example}[有理数域的完备化]\label{ex:completion-Q}
考虑度量空间 $(\mathbb{Q},d)$，其中
\[
d(p,q)=|p-q|,\qquad p,q\in\mathbb{Q}.
\]
我们用柯西列等价类来构造它的完备化。

\begin{enumerate}
  \item \textbf{第一步：取所有有理数柯西列的集合。}

  记
  \[
    C := \bigl\{ (x_n)_{n\in\mathbb{N}} \mid x_n\in\mathbb{Q},\ (x_n)\ \text{在}\ \mathbb{Q}\ \text{中是柯西列} \bigr\}.
  \]
  即 $C$ 中的每个元素都是一个有理数柯西序列。
  
  \item \textbf{第二步：用“差趋于 0”定义等价关系。}

  对任意两条柯西列 $(x_n),(y_n)\in C$，定义
  \[
    (x_n)\sim (y_n)
    \quad\Longleftrightarrow\quad
    \lim_{n\to\infty}|x_n-y_n|=0.
  \]
  即：两条柯西列的项差趋于 $0$ 时，将它们视为“表示同一个极限点”的不同方式，从而视为等价。
  
  记等价类集合为
  \[
    \widehat{\mathbb{Q}} := C/\sim.
  \]

  \item \textbf{第三步：在等价类上定义加法、数乘和“距离”。}

  对等价类 $[x_n],[y_n]\in\widehat{\mathbb{Q}}$，定义
  \[
    [x_n]+[y_n] := [x_n+y_n],\qquad
    \lambda\cdot[x_n]:=[\lambda x_n],\quad \lambda\in\mathbb{Q},
  \]
  以及
  \[
    \widehat d\bigl([x_n],[y_n]\bigr)
    :=\lim_{n\to\infty} |x_n-y_n|.
  \]
  可以验证上述定义与代表元的选取无关，且 $\widehat d$ 是 $\widehat{\mathbb{Q}}$ 上的度量，从而 $(\widehat{\mathbb{Q}},\widehat d)$ 成为一个度量线性空间。

  \item \textbf{第四步：把 $\mathbb{Q}$ 等距嵌入到 $\widehat{\mathbb{Q}}$ 中。}

  定义映射
  \[
    \iota:\mathbb{Q}\longrightarrow\widehat{\mathbb{Q}},\qquad
    q\longmapsto [(q,q,q,\dots)],
  \]
  即每个 $q\in\mathbb{Q}$ 对应常值列 $(q,q,q,\dots)$ 的等价类。显然 $\iota$ 是单射，并且
  \[
    \widehat d\bigl(\iota(p),\iota(q)\bigr)
    =\lim_{n\to\infty}|p-q|
    =|p-q|=d(p,q),
  \]
  因此 $\iota$ 是一个等距嵌入，$\mathbb{Q}$ 可以视为 $\widehat{\mathbb{Q}}$ 的一个稠密子集。

  \item \textbf{第五步：验证完备性与与 $\mathbb{R}$ 的同构性（思想说明）。}

  一方面，可以证明 $(\widehat{\mathbb{Q}},\widehat d)$ 是完备的：在 $\widehat{\mathbb{Q}}$ 中任取一条柯西列 $([x_n^{(k)}])_{k}$，可构造出一条在 $\mathbb{Q}$ 中的“二重序列”，再适当对角化，得到 $C$ 中的一条柯西列，其等价类作为极限点，从而说明 $\widehat{\mathbb{Q}}$ 完备。

  另一方面，每个实数 $r\in\mathbb{R}$ 都可以由某条有理数柯西列 $(q_n)$ 收敛到它（即 $q_n\to r$），而不同的实数给出不同的等价类，从而得到一个保持加法、乘法与距离的双射
  \[
    \mathbb{R}\ \cong\ \widehat{\mathbb{Q}}.
  \]
  因此，$\widehat{\mathbb{Q}}$ 就是通常意义下的实数轴 $\mathbb{R}$。
\end{enumerate}
综上，$(\widehat{\mathbb{Q}},\widehat d)$ 是 $(\mathbb{Q},|\cdot|)$ 的完备化，并与 $(\mathbb{R},|\cdot|)$ 等距同构。
\end{example}
\section{压缩映射原理}
考虑常微分方程
\[
\begin{cases}
\dfrac{dx}{dt} = f(x,t), \\[6pt]
x\big|_{t=t_0} = x_0.
\end{cases}
\]

这个问题等价于求解积分方程
\[
x(t) = x_0 + \int_{t_0}^t f(x(\tau), \tau) \, d\tau.
\]

如果设算子 $T$ 为
\[
(Tx)(t) = x_0 + \int_{t_0}^t f(x(\tau), \tau) \, d\tau,
\]
则解这个积分方程的问题，就等价于在某个函数空间中求解不动点方程
\[
Tx = x.
\]

因此，研究完备度量空间中算子映射的不动点问题，是解决微分方程的重要途径。接下来我们将重点研究压缩映射的不动点定理。
\begin{theorem}[压缩映射原理]
设 $(X,d)$ 是完备度量空间，$T:X \to X$，并且存在常数 $0<\theta<1$，使得对任意 $x,y \in X$，有不等式
\[
d(Tx,Ty) \leq \theta d(x,y).
\]
则存在唯一的 $\bar{x} \in X$，使得
\[
T\bar{x} = \bar{x}.
\]
\end{theorem}
解方程kx=0
两边都加一个x，然后前面kx+x=T(x)
\begin{dxtips}[压缩映射原理总结]
\begin{enumerate}

  \item 也就是它总能稳定地缩小任意两点的距离。
  \item 结论：在完备空间里，这样的 $T$ 必然有一个唯一的不动点 $\bar{x}$，并且不断迭代 $T$（即 $x, T(x), T^2(x), \dots$）会收敛到 $\bar{x}$。
\end{enumerate}

\end{dxtips}
\begin{example}[实数上的压缩映射]
在度量空间 $(\mathbb{R},|\cdot|)$ 上令 $T(x)=\frac{x+3}{4}$。
则 $\forall x,y$ 有 $|T(x)-T(y)|\le \tfrac14|x-y|$，故 $T$ 为压缩映射。
其唯一不动点为解 $x=T(x)$ 得 $x=1$。
任意初值 $x_0$ 的迭代满足
\[
x_{n}=1+\Big(\tfrac14\Big)^{n}(x_0-1)\xrightarrow[n\to\infty]{}1.
\]
\end{example}
\href{https://www.bilibili.com/video/BV1tJpizREED?vd_source=aa3cbf11c6c44a57861db6554a266722}{\textbf{稳定点与压缩映射定理} 视频链接}
\begin{theorem}[压缩映射推广]
设 $(X,d)$ 是完备距离空间，$T:X \to X$ 是一个映射。  
如果存在自然数 $n_0$ 和常数 $0 \leq \theta < 1$，使得对所有 $x,y \in X$ 都有
\[
d\!\left(T^{n_0}x,\,T^{n_0}y\right) \leq \theta\, d(x,y),
\]
则 $T$ 在 $X$ 中有唯一的不动点。
\end{theorem}
	•	条件说：虽然 T 本身可能不是压缩映射，但它的某个迭代 $T^{n_0}$ 是压缩映射（收缩映射）。
	•	因为压缩映射在完备空间里有唯一不动点，所以这里推广到 T 也有唯一不动点。

\begin{proposition}[不动点的定义]
给定一个映射 $T:X \to X$，如果某个点 $\bar{x} \in X$ 在映射下保持不变：
\[
T(\bar{x}) = \bar{x},
\]
那么称 $\bar{x}$ 是 $T$ 的一个\textbf{不动点}（fixed point）。  

换句话说，这个点被映射“固定住了”，不会随着变换移动，所以叫 fixed point（不动点）。
\end{proposition}
\begin{dxtips}
    因为迭代都是一维度的映射可以将y引入，其实求解就是求交点，y=x是一条线，然后y=T（x)是另外一条线把x输入也就从x做一条垂线与T（x)的交点。
\end{dxtips}
在实数轴上一维情况下，迭代映射
\[
x_{n+1} = T(x_n)
\]
本来只是点在数轴上来回跳动，我们无法直观看到它是怎样逐步逼近不动点的。

于是引入二维坐标系，把原来的数轴作为横坐标 $x$，再额外引入一个纵坐标 $y$：
\begin{enumerate}
  \item 画出函数图像 $y = T(x)$，以及对角线 $y = x$：
  \begin{itemize}
    \item 对角线 $y = x$ 表示“不动点条件” $T(x) = x$；
    \item 它和曲线的交点就是不动点。
  \end{itemize}

  \item 把原来的一维迭代嵌入二维几何：
  \begin{itemize}
    \item 从点 $(x_n, x_n)$ 出发，竖直移动到 $(x_n, T(x_n))$（表示计算函数值）；
    \item 再水平移动到 $(T(x_n), T(x_n))$（表示更新下一步迭代）。
  \end{itemize}
\end{enumerate}

这样，原本抽象的数列 $\{x_n\}$ 就变成了一个在平面上“折线爬行”的轨迹。  
这个轨迹逐步螺旋/阶梯式地靠近交点 $(\bar{x}, \bar{x})$，我们肉眼就能看到“压缩映射 $\to$ 收敛到唯一不动点”的过程。
\begin{dxtips}
    格基也是通过携带函数，然后最短向量就得到解。
\end{dxtips}
\\
\begin{remark}[直观解释：饮料比喻]
设 $A$ 表示红杯子（红色饮料，对应原始数列 $x_n$），
$B$ 表示蓝杯子（蓝色饮料，对应迭代映射 $T(x)$），
$C$ 表示空杯子（额外引入的维度 $y$）。

\begin{itemize}
  \item 如果直接把 $A$ 与 $B$ 混在一起，就会乱成一团，难以分清过程。
  \item 更好的方式是引入一个空杯子 $C$：
    \begin{itemize}
      \item 倒 $A \to C$：先把红饮料放入空杯子（横坐标 $x$）；
      \item 倒 $B \to A$：再把蓝饮料作用到红饮料（映射 $T(x)$）；
      \item 倒 $C \to B$：空杯子承载了两者的交替作用（纵坐标 $y = T(x)$）。
    \end{itemize}
\end{itemize}

这样，$C$ 这个“空杯子”并不改变原有饮料（数学问题），
却提供了一个新的空间，把原本混杂的信息分离、显性化，
从而顺便看清迭代逐步收敛到共同点（不动点）的过程。
\end{remark}

\section{拓扑空间}
\begin{definition}[拓扑所有开集的集合]定义[ 1.5.1]
设 $X$ 是任一集，$\tau$ 是 $X$ 的子集构成的集族，且满足条件：
\begin{enumerate}
  \item 集 $X$ 与空集 $\varnothing$ 属于 $\tau$；
  \item $\tau$ 中任意个集的并集属于 $\tau$；
  \item $\tau$ 中任意有限个集的交集属于 $\tau$。
\end{enumerate}
则称 $\tau$ 是 $X$ 上的一个拓扑。集 $X$ 上定义了拓扑 $\tau$，称它是一个拓扑空间，记为 $(X,\tau)$。
在不致引起混淆时简记为 $X$。凡属于 $\tau$ 中的集称为开集。
\end{definition}
\begin{dxtips}
拓扑的定义就是：在一个集合 X 上，挑出一族子集 \tau，并且让它满足“包含空集和全集、任意并、有限交仍在里面”这三个条件。
•	这族子集 \tau 就被叫做 拓扑，而其中的集合被叫做 开集。

所以，拓扑这个动作，本质就是人为规定哪些子集要被当作开集。
不同的规定方式，就得到了不同的拓扑空间。
\end{dxtips}
\begin{example}[不同的拓扑]
设 $X=\mathbb{R}$。

\begin{enumerate}
  \item \textbf{通常拓扑（standard topology）}  
  拓扑基取所有开区间 $(a,b)$。因此开集就是任意开区间以及它们的并。  
  例如：
  \[
  (0,1), \quad (-\infty,0)\cup(1,\infty), \quad \mathbb{R}
  \]
  都是开集，而 $[0,1]$ 不是开集。

  \item \textbf{离散拓扑（discrete topology）}  
  定义拓扑 $\tau=\mathcal{P}(\mathbb{R})$，即所有 $\mathbb{R}$ 的子集都是开集。  
  例如：
  \[
  \{0\}, \quad \{1,2,3\}, \quad \mathbb{Q}
  \]
  都是开集。  
  \textbf{补充说明：} 离散拓扑是最“细”的拓扑（finest topology），因为它把所有子集都视为开集。它有一个重要性质：\emph{在离散拓扑下，任意函数 $f:X\to Y$ 都是连续的}，因为任意子集的原像仍然是开集。

  \item \textbf{平凡拓扑（trivial topology）}  
  定义拓扑 $\tau=\{\varnothing,\mathbb{R}\}$，即只有空集和全集是开集。  
  在这个拓扑下，$(0,1)$ 不是开集。

\end{enumerate}
\end{example}
\begin{itemize}
  \item 同一集合可以引进不同的拓扑，成为不同的拓扑空间。
  \item 设同一集合 $X$ 上有两个拓扑 $\tau_1,\ \tau_2$。如果 $\tau_2\subset \tau_1$
  （真包含可写作 $\tau_2\subsetneq \tau_1$），则称拓扑 $\tau_2$ 比拓扑 $\tau_1$ 弱，
  或者称拓扑 $\tau_1$ 比拓扑 $\tau_2$ 强。
\end{itemize}
\begin{proposition}[诱导拓扑与基本概念]
设 $(X,\tau)$ 是拓扑空间，$A\subset X$。令
\[
\tau_A = \{ A\cap G \mid G\in \tau \},
\]
则称 $\tau_A$ 为 $\tau$ 在 $A$ 上的诱导拓扑，称 $(A,\tau_A)$ 为 $(X,\tau)$ 的子空间。

\begin{itemize}
  \item \textbf{闭集：} 开集的余集。
  \item \textbf{邻域：} 若开集 $G$ 含有 $x$，则称 $G$ 为 $x$ 的邻域。
  \item \textbf{接触点：} 若 $x$ 的每个邻域都包含 $E$ 中的点，则称 $x$ 为 $E$ 的接触点。
  \item \textbf{极限点：} 若 $x$ 的每个邻域都包含 $E$ 中不同于 $x$ 的点，则称 $x$ 为 $E$ 的极限点。
  \item \textbf{$E$ 的闭包：} $E$ 的接触点全体。
  \item \textbf{$E$ 的内部：} $E$ 中所有开集的并集。
\end{itemize}
\end{proposition}
我们原来有一个大空间 $(X,\tau)$，上面已经定义好了拓扑。
现在如果只关心某个子集 $A\subset X$，那 $A$ 上也要有拓扑，怎么办？

\medskip
\noindent
\textbf{诱导拓扑} 的定义就是告诉我们：子集 $A$ 的拓扑要“继承”自 $X$ 的拓扑。

\medskip
\noindent
具体做法：
\begin{itemize}
  \item $A$ 上的开集定义为
  \[
  \tau_A=\{\,A\cap G \mid G\in\tau\,\}.
  \]
  也就是说：大空间里的开集 $G$ 与 $A$ 相交，得到的就是小空间中的开集。
\end{itemize}
\begin{definition}[拓扑基]
设 $(X,\tau)$ 是一个拓扑空间，$\mathcal{B}$ 是 $\tau$ 的子集族，
使得 $X$ 中的每一开集可表为 $\mathcal{B}$ 中集的并集，称 $\mathcal{B}$ 是 $(X,\tau)$ 的\textbf{拓扑基}。
这样，我们可以指定空间的拓扑基来给出拓扑。
\end{definition}
拓扑基就是能够生成整个拓扑的集合
\begin{theorem}[拓扑基的性质]
设 $X$ 是任一集合，$\mathcal{B}$ 是 $X$ 的子集构成的集族，
则 $\mathcal{B}$ 是 $X$ 上某一拓扑基，当且仅当 $\mathcal{B}$ 具有以下性质：
\begin{enumerate}
  \item 任一点 $x\in X$，存在 $G\in\mathcal{B}$，使得 $x\in G$；
  \item 如果 $x\in G_{1}\cap G_{2}$，其中 $G_{1},G_{2}\in\mathcal{B}$，则存在 $G_{3}\in\mathcal{B}$，
  使得 $x\in G_{3}\subset G_{1}\cap G_{2}$。
\end{enumerate}
\end{theorem}
\begin{theorem}[拓扑基的判定]
设族 $\mathcal{B} \subset \tau$，则 $\mathcal{B}$ 是给定拓扑 $\tau$ 的基，当且仅当满足以下条件：  
对任意 $G \in \tau$ 及每一点 $x \in G$，存在 $G_x \in \mathcal{B}$，使得
\[
x \in G_x \subset G.
\]
\end{theorem}
\begin{example}
 距离空间中，所有开球的全体构成拓扑基。
\end{example}
\begin{example}
\[
B=\{(q_1,q_2)\mid \forall q_i\in \mathbb{Q},\, i=1,2\}
\]
构成欧式空间 $\mathbb{R}$ 的拓扑基。
\end{example}
\noindent\textcolor{red}{\large\textbullet}\quad 
如果拓扑空间中存在一个由至多可数个集构成的拓扑基，称这个空间具有\textbf{可数基}（或满足\textbf{第二可数公理}）。
\noindent\textcolor{red}{\large\textbullet}\quad 
具有可数基的距离空间是\textbf{可分}的。

\begin{definition}[连续映射]
设 $X, Y$ 是两个拓扑空间，$f : X \to Y$ 是一个映射，$x_0 \in X$。  
如果对点 $y_0 = f(x_0)$ 的任意邻域 $U_{y_0}$，存在点 $x_0$ 的邻域 $V_{x_0}$，使得
\[
f(V_{x_0}) \subset U_{y_0},
\]
则称 $f$ 在点 $x_0$ 连续；如果 $f$ 在每一点 $x \in X$ 连续，则称 $f$ 是一个连续映射。  

拓扑空间到数直线上（$\mathbb{R}$ 上）的连续映射称为连续函数。
\end{definition}
\begin{theorem}[连续映射的等价条件]
设 $X, Y$ 是拓扑空间，$f: X \to Y$ 是一个映射。则以下三个命题等价：
\begin{enumerate}
  \item $f$ 在 $X$ 上是连续的；
  \item 对 $Y$ 中的任意开集 $G$，$f^{-1}(G)$ 是 $X$ 中的开集；
  \item 对 $Y$ 中的任意闭集 $F$，$f^{-1}(F)$ 是 $X$ 中的闭集。
\end{enumerate}
\end{theorem}
\begin{proposition}[开映射与同胚]
设 $X, Y$ 是拓扑空间，$f: X \to Y$ 为一个映射。

\begin{enumerate}
  \item 若 $f$ 将 $X$ 中的任意开集映为 $Y$ 中的开集，则称 $f$ 为\textbf{开映射}；
  \item 若 $f$ 为双射，且 $f$ 与其逆映射 $f^{-1}$ 都连续，则称 $f$ 为\textbf{同胚映射}；
  \item 若两个拓扑空间之间存在同胚映射，则称这两个空间\textbf{同胚}。
\end{enumerate}
\end{proposition}


\begin{theorem}[定理 1.5.5]
设 $X,Y,Z$ 是拓扑空间，$f:X \to Y$ 与 $g:Y \to Z$ 是连续映射，则
\[
h(x) = g(f(x)), \qquad x \in X
\]
是空间 $X$ 上到空间 $Z$ 中的连续映射。
\end{theorem}
\subsection{分离公理}
\begin{definition}[$T_1$ 分离公理]
设 $X$ 为拓扑空间。若任取 $x,y\in X$ 且 $x\neq y$，存在 $x$ 的邻域 $U_x$ 不含 $y$，且存在 $y$ 的邻域 $U_y$ 不含 $x$，则称 $X$ 满足 \textbf{$T_1$ 分离公理}。
\end{definition}
也就是x和y是可以分开的，判定就是每一个人都有一个领域不包含对方。
\begin{definition}[$T_2$ 分离公理（Hausdorff 公理）]
设 $X$ 为拓扑空间。若任取 $x,y\in X$ 且 $x\neq y$，存在分别属于 $x$ 和 $y$ 的邻域 $G_x,G_y$，使得 $G_x\cap G_y=\varnothing$，则称 $X$ 满足 \textbf{$T_2$ 分离公理}，或称 $X$ 为 \textbf{Hausdorff 空间}。
\end{definition}

\begin{definition}[$T_3$ 分离公理]
设 $X$ 为 $T_1$ 空间。若对任意 $x\in X$ 及不含 $x$ 的闭集 $A$，存在开集 $U_x,U_A$，满足 $x\in U_x$，$A\subset U_A$ 且 $U_x\cap U_A=\varnothing$，则称 $X$ 满足 \textbf{$T_3$ 分离公理}。
\end{definition}

\begin{definition}[$T_4$ 分离公理]
设 $X$ 为 $T_1$ 空间。若对任意两个不相交闭集 $A,B\subset X$，存在开集 $U_A,U_B$，满足 $A\subset U_A$，$B\subset U_B$ 且 $U_A\cap U_B=\varnothing$，则称 $X$ 满足 \textbf{$T_4$ 分离公理}，或称 $X$ 为 \textbf{正规空间}。
\end{definition}

\noindent
前面的例 1.5.5 是一个不满足 $T_1$ 分离公理的拓扑空间的例子。
\begin{theorem}[定理 1.5.6]
拓扑空间 $X$ 是 $T_1$ 的，当且仅当 $X$ 中任一单点集是闭集。
\end{theorem}
\begin{theorem}[定理 1.5.6]
拓扑空间 $X$ 是 $T_1$ 的，当且仅当 $X$ 中任一单点集是闭集。
\end{theorem}
\subsection{紧性}
\begin{definition}[紧空间]
设 $X$ 是拓扑空间。若存在开集族 $\{G_\alpha \subset X \mid \alpha \in I\}$ 满足
\[
X \subset \bigcup_{\alpha \in I} G_\alpha,
\]
称 $\{G_\alpha \mid \alpha \in I\}$ 为 $X$ 的一个\textbf{开覆盖}。  
若对 $X$ 的任意开覆盖都存在有限子覆盖，即存在有限指标集
\[
\{\alpha_1,\alpha_2,\dots,\alpha_n\}\subset I,
\quad 使得\quad
X \subset \bigcup_{k=1}^n G_{\alpha_k},
\]
则称 $X$ 是\textbf{紧的}（或称 $X$ 为\textbf{紧空间}）。
\end{definition}
任意开覆盖都得有有限子覆盖
\begin{definition}[紧集]
设 $A$ 是拓扑空间 $X$ 的子集。  
若对 $A$ 的任意开覆盖 $\{G_\alpha \mid \alpha \in I\}$（其中 $G_\alpha$ 为 $X$ 中的开集）  
都存在有限子覆盖，则称 $A$ 为 $X$ 的\textbf{紧子集}。  

等价地说：对 $X$ 中任意开集族 $\{G_\alpha \subset X \mid \alpha \in I\}$，若其覆盖 $A$，  
则必存在有限个 $G_{\alpha_1},\dots,G_{\alpha_n}$ 覆盖 $A$。
\end{definition}
\begin{theorem}[紧空间的基本性质]
设 $X$ 为拓扑空间，则紧空间具有以下性质：
\begin{enumerate}
  \item $X$ 是紧的，当且仅当 $X$ 具有\textbf{有限交性质}；
  \item 紧空间中的任意闭子集是紧的；
  \item Hausdorff 空间中的每一个紧子集是闭的；
  \item （Tychonoff 定理）任意多个紧空间的乘积空间是紧的。
\end{enumerate}

\end{theorem}

\begin{theorem}[紧空间的连续映射性质]
设 $f:X_1\to X_2$ 是拓扑空间之间的连续映射，则有以下性质：

\begin{enumerate}
\item 若 $X_1$ 是紧空间，则 $f(X_1)$ 是 $X_2$ 中的紧集；
  
\item 若 $f:X\to \mathbb{R}$ 在紧空间 $X$ 上连续，则 $f$ 必有最大值与最小值；
  
\item 若 $X$ 是紧 Hausdorff 空间，且 $f:X\to Y$ 是到 Hausdorff 空间 $Y$ 的一一连续映射，则 $f$ 是同胚映射。
\end{enumerate}
\end{theorem}

\begin{definition}[列紧]
设 $A$ 是距离空间 $X$ 的子集。  
如果 $A$ 中任意点列必包含一个在 $X$ 中收敛的子列，  
则称 $A$ 是\textbf{列紧集}。  
如果 $X$ 是列紧集，则称 $X$ 是\textbf{列紧空间}。
\end{definition}

\begin{remark}
\begin{itemize}
  \item 列紧集的子集是列紧的；
  \item 列紧空间是完备的。
\end{itemize}
\end{remark}
列紧比完备强，完备需要柯西列收敛。
\begin{definition}[有界集]
称距离空间 $X$ 的任意子集 $A$ \textbf{有界}，  
若存在开球 $S(x_0, r)$ 使得
$
A \subset S(x_0, r).
$
\end{definition}

\begin{remark}
在欧氏空间 $\mathbb{R}$ 中，集合 $A$ 有界等价于 $A$ 列紧。  
但在一般的距离空间中，集合 $A$ 有界不能推出 $A$ 列紧。
\end{remark}

\begin{definition}[全有界集totally bounded space]
设 $A,B$ 是距离空间 $X$ 的子集，$\varepsilon>0$ 是一个给定的正数。  
如果对每个 $x\in A$，都存在 $y\in B$，使得
\[
d(x,y)<\varepsilon,
\]
则称 $B$ 是 $A$ 的一个 \textbf{$\varepsilon$-网}。

若 $A\subset X$，并且对任意 $\varepsilon>0$，$A$ 都有有限 $\varepsilon$-网，  
则称 $A$ 是 \textbf{全有界集}。
\end{definition}
有限个，$\epsilon$小了B的个数就要增加才能覆盖但是不能超过有限个。无界没法有限覆盖。无界一定不全有界。

\begin{proposition}[全有界集的性质]
\begin{enumerate}[(1)]
  \item 全有界集一定是有界集；
  \item 全有界集 $A$ 的有限 $\varepsilon$-网可以取为 $A$ 的子集；
  \item 全有界集是可分的；进一步地，若距离空间 $X$ 全有界，则 $X$ 具有可数基。
\end{enumerate}
\end{proposition}
\begin{theorem}
\begin{enumerate}[(1)]
  \item 设 $A$ 是距离空间 $X$ 中的列紧集，则 $A$ 是全有界集；
  \item 如果 $X$ 是完备距离空间，则当 $A$ 是全有界集时，$A$ 是列紧集。
\end{enumerate}

在空间 $\mathbb{R}^n$ 中，有界集是全有界集；  
但在一般的距离空间中，此结论不一定成立。
\end{theorem}
把A盖上就能把点列盖上，覆盖是有限的点列是无限的所以有一个邻域里面有无限个点。有限维有界点列必有收敛子列，距离空间不一定了。L2，对角线元？证明完备以后收敛只用柯西列。抽取对角线
\begin{theorem}[列近闭集的子集是紧集]
设 $A$ 是距离空间 $X$ 的子集，则 $A$ 是紧集，当且仅当 $A$ 是列紧闭集。
\end{theorem}
\begin{remark}
 “列紧（sequentially compact）” 这一名称正是来源于它的定义性质：

\begin{itemize}
  \item “列 (sequence)” 对应序列；
  \item “紧 (compact)” 对应“极限不跑出集合”；
\end{itemize}

因此，\textbf{“列紧”} 的字面意思就是：

\text{“序列意义下的紧集” (a compact set in the sense of sequences).}

\end{remark}
\begin{theorem}[Arzelà 定理]
空间 $C[a,b]$ 的子集 $A$ 是列紧的，当且仅当以下两个条件同时成立：

\begin{enumerate}[(1)]
  \item $A$ 一致有界：存在常数 $K>0$，使得对任意 $x\in A$、$t\in[a,b]$，都有
  \[
  |x(t)| \le K;
  \]
  
  \item $A$ 等度连续：对任意 $\varepsilon>0$，存在 $\delta>0$，使得对任意 $x\in A$ 及任意 $t_1,t_2\in[a,b]$，若 $|t_1-t_2|<\delta$，则有
  \[
  |x(t_1)-x(t_2)|<\varepsilon.
  \]
\end{enumerate}
\end{theorem}
\begin{example}
记
\[
C[a,b] = \{\,x:[a,b]\to\mathbb{R} \mid x(t)\ \text{在}\ [a,b]\ \text{上连续}\,\},
\]
即所有定义在区间 $[a,b]$ 上的连续实值函数的集合。

\begin{itemize}
  \item $C[a,b]$：  
  例如 $x(t)=\sin t,\ x(t)=t^2,\ x(t)=e^t$ 都属于 $C[0,1]$。
  \item $A \subset C[a,b]$：表示从这些连续函数中挑出一部分（一个“函数族”）。  
  例如 $A=\{t, t^2, t^3, \dots\}$。
\end{itemize}

取
\[
A = \{\,x_n(t)=t^n : n=1,2,3,\dots\,\} \subset C[0,1].
\]

对任意 $x\in A$（例如 $x(t)=t^3$），以及任意 $t\in[0,1]$，都有
\[
|x(t)| = |t^3| \le 1,
\]
因此 $A$ 是\textbf{一致有界}的（可取 $K=1$）。

又因为对任意 $\varepsilon>0$，当 $|t_1-t_2|$ 很小时，$|t_1^n - t_2^n|$ 也会很小，
所以 $A$ 也是\textbf{等度连续}的。

由 Arzelà–Ascoli 定理可知，$A$ 在 $C[0,1]$ 中是\textbf{列紧}的。
\end{example}
\begin{proposition}[Arzelà--Ascoli 定理的独特价值]
Arzelà--Ascoli 定理是 Heine--Borel 定理在函数空间 $C[a,b]$ 中的推广，
它刻画了在无限维空间中哪些函数族是\textbf{相对紧}的——
即任意函数列都能抽出一致收敛子列的集合。

\begin{itemize}
  \item[\textbf{(1)}] \textbf{有限维情形：}  
  在 $\mathbb{R}^n$ 中，有界且闭的集合为紧集（Heine--Borel 定理）：
  \[
  \text{有界且闭} \;\Rightarrow\; \text{紧集.}
  \]

  \item[\textbf{(2)}] \textbf{函数空间情形：}  
  在无限维的函数空间 $C[a,b]$ 中，“有界闭”已不足以保证紧性。
  Arzelà--Ascoli 定理指出，若一族函数一致有界且等度连续，则它是列紧的：
  \[
  \text{一致有界且等度连续} \;\Rightarrow\; \text{列紧.}
  \]

  \item[\textbf{(3)}] \textbf{理论意义：}  
  Arzelà--Ascoli 定理揭示了函数空间中“收敛性”与“紧性”的深层联系，
  将数学分析中“一致收敛”的概念提升为拓扑意义下的“紧性”刻画，
  从而在无限维空间中建立了“收敛—结构”之间的桥梁。
\end{itemize}
\end{proposition}































\section{第一章 习题}

\begin{exercise}
设 $D$ 是 $[0,1]$ 区间上具有连续导数（在端点 $t=1,t=0$ 分别具有左、右导数）的实函数全体，在 $D$ 上定义
\[
d(x,y) = \sup_{0 \leq t \leq 1} |x(t) - y(t)| \;+\; \sup_{0 \leq t \leq 1} |x'(t) - y'(t)|.
\]
\begin{enumerate}
  \item 证明 $D$ 是距离空间；
  \item 指出 $D$ 中点列按距离收敛的意义；
  \item 证明 $D$ 是完备的。
\end{enumerate}
\end{exercise}

\begin{exercise}
证明如果 $d$ 是集 $X$ 上的距离，则
\[
d_1 = \frac{d}{1+d}
\]
也是 $X$ 上的距离。
\end{exercise}
\begin{exercise}
设 $d_1,d_2,\dots,d_m,\dots$ 是集 $X$ 上的距离，证明：

\begin{enumerate}
  \item $d = \sup\limits_{1 \leq i \leq m} d_i$；
  \item $d = \sqrt{d_1^2 + d_2^2 + \cdots + d_m^2}$；
  \item $d = \sum\limits_{k=1}^{\infty} \frac{1}{2^k} \cdot \frac{d_k}{1+d_k}$；
\end{enumerate}

中的每一个也是 $X$ 上的距离。
\end{exercise}

\begin{exercise}
设 $X$ 是在 $|z|<1$ 中解析且在 $|z|\leq 1$ 上连续的复函数的全体，在其中定义
\[
d(x,y) = \max_{|t|=1} |x(t) - y(t)|,
\]
证明 $(X,d)$ 是距离空间。
\end{exercise}

\begin{exercise}
在距离空间中，一个半径为 $4$ 的开球能否成为一个半径为 $3$ 的开球的真子集？
\end{exercise}

\begin{exercise}
证明在距离空间中，如果一个半径为 $7$ 的开球包含在一个半径为 $3$ 的开球中，则这两个球重合。
\end{exercise}

\begin{exercise}
证明在空间 $S$ 中，按距离收敛等价于按坐标收敛。
\end{exercise}

\begin{exercise}
试举一例说明有界集不是全有界集。
\end{exercise}

\begin{exercise}
证明距离空间中每一个 Cauchy 列是有界集。
\end{exercise}

\begin{exercise}
证明距离空间的完备子空间是闭子空间。
\end{exercise}

\begin{exercise}
证明如果距离空间是可分的，则它的任意子空间也是可分的；反之，如果距离空间不可分，它的子空间是否也不可分？
\end{exercise}
\begin{exercise}
设 $(X,d)$ 是距离空间，$A \subset X$，令
\[
f(x) = \inf_{y \in A} d(x,y), \quad (x \in X).
\]
证明 $f(x)$ 是 $X$ 上的连续函数。
\end{exercise}

\begin{exercise}
设 $F_1, F_2$ 是距离空间 $X$ 中不相交的闭集，证明存在 $X$ 上的连续函数 $f(x)$，使得当 $x \in F_1$ 时 $f(x)=0$，当 $x \in F_2$ 时 $f(x)=1$。
\end{exercise}

\begin{exercise}
设 $X$ 是距离空间，证明：如果在 $X$ 中，任一半径趋于零的闭球套具有非空交，则空间 $X$ 是完备的。
\end{exercise}

\begin{exercise}
设 $X$ 是完备距离空间，$\mathcal{F}$ 是 $X$ 上的实连续函数族且具有性质：对每一 $x \in X$，存在常数 $M_x > 0$，使得对每一 $F \in \mathcal{F}$，
\[
|F(x)| \leq M_x.
\]
证明：存在开集 $U$ 及常数 $M>0$，使得对每一 $x \in U$ 及所有 $F \in \mathcal{F}$，
\[
|F(x)| \leq M.
\]
\end{exercise}
\begin{exercise}
举例说明，在压缩映射原理中：
\begin{enumerate}
  \item 空间完备性条件不可少；
  \item 映射 $T$ 所满足的条件不能代之以条件：
  \[
  d(Tx,Ty) < d(x,y), \quad (x \neq y).
  \]
\end{enumerate}
\end{exercise}

\begin{exercise}
证明：存在闭区间 $[0,1]$ 上的连续函数 $x(t)$，使得
\[
x(t) = \tfrac{1}{2}\sin x(t) - a(t),
\]
其中 $a(t)$ 是给定的 $[0,1]$ 上的连续函数。
\end{exercise}

\begin{exercise}
设 $X$ 是完备距离空间，$T$ 是 $X$ 上到自身的映射，在闭球
\[
B = \{ x \in X : d(x_0,x) \leq r \}
\]
上，若 $d(Tx,Ty) \leq \theta d(x,y)$ 且 $d(x_0,Tx_0) < (1-\theta)r$，其中 $0 \leq \theta < 1$，证明 $T$ 在 $B$ 上有唯一不动点。
\end{exercise}

\begin{exercise}
设 $(t_0,s_0)\in \mathbb{R}^2$，$f(t,s)$ 在 $(t_0,s_0)$ 的邻域 $N$ 中连续，$s_0=f(t_0,s_0)$，$f'_s(t,s)$ 在 $N$ 中存在且在 $(t_0,s_0)$ 连续，并且 $f'_s(t_0,s_0)=0$。用压缩映射原理证明：存在 $\delta > 0$，$x(t)\in C[t_0-\delta,\,t_0+\delta]$，使得
\[
s_0 = x(t_0), \quad x(t) = f(t,x(t)), \quad t\in [t_0-\delta,\,t_0+\delta].
\]
\end{exercise}

\begin{exercise}
设 $X$ 是紧距离空间，$T$ 是 $X$ 上到自身的映射且满足条件：对任意 $x,y \in X$，当 $x\neq y$ 时，
\[
d(Tx,Ty) < d(x,y).
\]
证明 $T$ 在 $X$ 上有唯一不动点。
\end{exercise}
\begin{exercise}
设 $X=\{a,b\}$，令 $\tau=\{X,\varnothing,\{b\}\}$。证明：
(1) $\tau$ 是 $X$ 上的一个拓扑；
(2) 集合 $\{a\}$ 是闭集；
(3) $\overline{\{b\}}=X$。
\end{exercise}

\begin{exercise}
设 $X$ 为任一集合，$\mathcal{B}_1,\mathcal{B}_2$ 是由 $X$ 的子集构成的基，并据此分别确定 $X$ 上的拓扑 $\tau_1,\tau_2$。证明：$\tau_1\subset\tau_2$ 当且仅当对任意 $G_1\in\mathcal{B}_1$ 及任意点 $x\in G_1$，都存在 $G_2\in\mathcal{B}_2$ 使得 $x\in G_2\subset G_1$。
\end{exercise}

\begin{exercise}
举出一个拓扑空间的例子，它是 $T_2$（Hausdorff）的，但不是 $T_3$（正则 $T_1$）的。
\end{exercise}

\begin{exercise}
设 $X$ 为 $T_1$ 拓扑空间，$E\subset X$，证明：点 $x\in X$ 是 $E$ 的极限点，当且仅当 $x$ 的任意邻域都包含 $E$ 中无穷多个点。若不假设 $X$ 为 $T_1$，此结论是否仍成立？
\end{exercise}

\begin{exercise}
设拓扑空间 $(X,\tau)$ 满足第二可数公理。证明：从 $X$ 的任意开覆盖中可以选出由至多可数个开集构成的子覆盖。
\end{exercise}






































































